\section{Discussion}
\label{sec:discussion}

The planner experiment shows that the service split matches where the computing time goes. A* search and waypoint synthesis take only a small part of the planning time, while CFMS invocation and dense trajectory generation take most of it and also cause the dependence on route length. On the Hamburg network, graph routing is therefore not the main bottleneck; the slow part is turning a route into a dense 4D trajectory. The regressions explain why: the number of trajectory samples grows almost linearly with route length, and trajectory-generation time grows linearly with the number of samples, so longer routes take longer mainly because they have more samples.

This matters for route updates and contingency planning. If a corridor edge has to be removed because of a closure, a geofence or a conflict, the short A* runtime means that a new lateral route can be found quickly. The larger delay comes from generating the new 4D trajectory and, with conflict checking, from checking it again against all stored peers.

The conflict experiment shows the NDMap and re-routing loop working with both sparse and dense traffic. Over the full day, most candidates were conflict-free at the first check, and the few conflicts were solved by removing edges for that request. With a shorter scheduling window, more flights overlap in time, there are clearly more conflicts, and we see outcomes that the sparse run cannot show: requests hitting the retry limit and requests running out of lateral alternatives. Mapping stalls did not occur in any run, so this outcome is implemented and logged but was never observed. The scripted encounters confirm the expected behavior: head-on and crossing encounters are detected and solved laterally, while merges at the same pad and landing time, and in-trail traffic on the same route, cannot be solved by lateral re-routing alone. All these outcomes are results of planning and deconfliction before submission, not USSP authorization decisions.

It is worth being clear about how time is handled, because overlap in time and a real conflict are easy to mix up. Two trajectories that share a scheduling window, or even their whole flight time, are not in conflict because of that. The detector needs one moment at which both positions break the horizontal and vertical minima at the same time. Per pair, overlap is rare in the full-day run (\cfOverlapPairs{} of \cfComparisons{} pairs). Per request it is more common, because a request has many peers and needs only one of them to overlap; \cfWithOverlap{} of \cfWithPeers{} requests have at least one overlapping peer, but only \cfConflicted{} of them have a conflict. The same holds in the shorter windows. So conflict and overlap numbers always have to be read together with what they are counted over.

The latency breakdown shows where the time goes with conflict checking. Planning stays at tens of milliseconds, but every conflict check loads the candidate and all peer trajectories from the trajectory store, builds the detector state in a new process, and only then runs the check. This cost grows with the number of peers and becomes much larger than trajectory generation; resolved requests pay it again for every attempt. We chose a stateless detector to keep the state consistent, and this cost is the price for it.

The most useful improvement is therefore a detector state that is kept or updated step by step, so that unchanged peer trajectories are not loaded and indexed again for every check. After that come caching in the trajectory store and a clearer split between outputs for visualization and inputs for conflict checking. Resolution should also be extended to time and altitude, because our encounter cases show that lateral re-routing alone cannot solve shared-destination or dense in-trail conflicts and can use up the alternatives in the graph. The single cruise layer makes this even clearer. With all cruise segments at the same altitude, the vertical criterion is always met between two cruising aircraft, so en-route separation depends only on horizontal distance and lateral re-routing is the only option the loop has. A corridor structure with several layers would make the vertical criterion useful en route and give the loop a second dimension to work with. The conflict numbers in this paper should be read with this single-layer assumption in mind.

We also want to look critically at the corridor assumption behind this work. U-space does not require corridors. U-space services are provided in a region or volume in which operators can plan their trajectories freely, and corridors are only one option that some operators want to use for early UAM operations, as studied for example in the SESAR project EUREKA~\cite{sesar_eureka_project}. We chose corridors because they map directly onto a graph, make procedure design and conflict geometry simpler, and are a plausible structure for first operations. On the other hand, airspace studies show that more structure usually means more predictability but less capacity and flexibility: layered or less constrained airspace can handle more traffic at high densities~\cite{sunil2015metropolis,schneider2014metropolis,abdellaoui2025_airspace_uam,schuchardt2023_atm_uam_review}. For our approach this trade-off is built in, since a corridor graph only offers a limited set of lateral alternatives to the re-routing loop---and the dense and scripted cases reach this limit. The architecture itself does not depend on corridors, since the fleet webserver could also provide other kinds of route networks, but the capacity-related results are specific to the corridor model we tested.

The limitations of Section~\ref{sec:limitations} still apply, and two of them limit how far the statistics go. The planner experiment has \paObs{} timing observations but only \paPairs{} different routes, so the route-length results rest on \paPairs{} independent points. The density runs use one fixed OD sequence, order and seed; they show how our loop behaves when the window gets shorter, but the shares we report belong to this one run and are not conflict probabilities in general. Within these limits, the results describe the computing behavior of our pipeline, not the capacity of an operational system.
